\documentclass[10pt]{article}
\usepackage[utf8]{inputenc}
\usepackage{geometry}
\geometry{a4paper, margin=0.75in}
\usepackage{amsmath}
\usepackage{amssymb}
\usepackage{booktabs}
\usepackage{hyperref}
\usepackage{enumitem}
\setlist{nosep}

\title{DeepLOB Implementation Report: Limit Order Book Prediction with Deep Learning}
\author{Deep Learning Project}
\date{\today}

\begin{document}

\maketitle

\begin{abstract}

DeepLOB is a hybrid deep learning architecture for limit order book (LOB) prediction, achieving $\sim$80\% F1 scores on the FI-2010 benchmark. By combining Convolutional Neural Networks (CNNs), Inception modules, and LSTMs, it effectively captures spatial LOB structure, multi-scale temporal patterns, and long-term dependencies. This report details the architecture, implementation, and results.

\end{abstract}

\section{Introduction and Data}

\subsection{The Challenge}
Limit Order Books contain high-dimensional, noisy, and non-stationary data (bid/ask prices and volumes across levels). DeepLOB aims to predict price direction (Up, Stationary, Down) from raw LOB snapshots without handcrafted features.

\subsection{FI-2010 Dataset}
The standard FI-2010 dataset consists of 4 million samples from 5 Finnish stocks.
\begin{itemize}
    \item \textbf{Input}: 100 historical snapshots $\times$ 40 features (10 levels $\times$ \{Price, Vol\} for Bid/Ask).
    \item \textbf{Normalization}: The dataset uses static z-score normalization. However, for real-world deployment, \textbf{rolling normalization} (e.g. 5-day window) is required.
    \item \textbf{Labels}: Derived from smoothed future returns. The recommended labeling method compares smoothed future prices to smoothed past prices to reduce noise:
    \[l_t = \frac{m^+(t) - m^-(t)}{m^-(t)}, \quad y_t = \text{threshold}(l_t)\]
\end{itemize}

\section{Architecture and Implementation Analysis}

DeepLOB processes data through three sequential stages, each addressing a specific domain challenge.

\subsection{Convolutional Block: Spatial Features \& Stride}
The first block extracts ``microstructure'' features. A key design choice is the \textbf{(1$\times$2) stride} in the first layer.
\begin{itemize}
    \item \textbf{Why}: Features are ordered $\{P_a, V_a, P_b, V_b\}$. A standard stride of 1 would correlate semantically unrelated pairs (e.g., $V_a$ and $P_b$).
    \item \textbf{Solution}: Stride 2 forces the kernel to process valid $(P, V)$ pairs only, respecting the LOB's semantic structure.
    \item \textbf{Structure}:
    \begin{enumerate}
        \item Conv(1$\times$2, stride 1$\times$2) $\rightarrow$ Aggregates Price/Vol.
        \item Conv(1$\times$2, stride 1$\times$2) $\rightarrow$ Approximates the \textbf{micro-price} $p_{\text{micro}} = I \cdot p_a + (1-I) \cdot p_b$ (where $I$ is imbalance).
        \item Conv(1$\times$10) $\rightarrow$ Global LOB summary.
    \end{enumerate}
\end{itemize}

\subsection{Inception Module: Multi-Scale Discovery}
Financial patterns occur at unknown timescales. Instead of selecting a fixed window, the Inception module applies parallel filters:
\begin{itemize}
    \item \textbf{Paths}: $1\times1$ (bottleneck), $3\times1$ (short-term), $5\times1$ (medium), $10\times1$ (long-term), and MaxPool.
    \item \textbf{Result}: The network automatically learns which timescales are predictive by concatenating these diverse feature maps (128 features total).
\end{itemize}

\subsection{LSTM: Temporal Dependencies}
An LSTM (64 units) processes the sequence of Inception features.
\begin{itemize}
    \item \textbf{Efficiency}: Using an LSTM instead of a fully connected layer reduces parameters by $\sim$10x (60k vs 800k), preventing overfitting on the relatively small FI-2010 dataset.
    \item \textbf{Regularization}: No Batch Normalization is used due to non-stationarity. Leaky ReLU ($\alpha=0.01$) and early stopping are preferred.
\end{itemize}

\section{Experimental Setup}

The loss function used the categorical cross-entropy with Adam optimizer ($lr=0.01$ and $\epsilon=1$) and batch size 32.

\subsection{Implementation Checklist}
To reproduce DeepLOB, ensure the following:
\begin{itemize}[label=$\square$]
    \item \textbf{Data}: Shape (N, 100, 40). Normalized (z-score), labels one-hot.
    \item \textbf{Conv1}: Kernel (1,2), \textbf{Stride (1,2)}. (Critical!)
    \item \textbf{Inception}: Concatenate outputs on the \textit{feature} axis.
    \item \textbf{LSTM}: 64 units, returns last state.
    \item \textbf{Training}: Adam ($lr=0.01$), Batch size 32, Leaky ReLU ($\alpha=0.01$).
    \item \textbf{Validation}: Use separate validation set for early stopping.
\end{itemize}

\subsection{Train/Test Splits}
Two setups are used:
\begin{itemize}
    \item \textbf{Setup 1}: Rolling window cross-validation (5 folds).
    \item \textbf{Setup 2}: Fixed split (first 7 days train, last 3 days test) for direct comparison with baselines.
\end{itemize}

\subsection{Evaluation Metrics}
\begin{itemize}
    \item \textbf{Accuracy}: Overall correctness.
    \item \textbf{Precision}: Class-wise correctness.
    \item \textbf{Recall}: Class-wise completeness.
    \item \textbf{F1 Score}: Macro-averaged across classes to handle class imbalance.
    \item \textbf{Trading Simulation}: Naive strategy to assess economic value.
\end{itemize}

\section{Experimental Results}

\subsection{Performance (FI-2010)}
DeepLOB outperforms standard baselines (SVM, MLP, CNN, LSTM) and the attention-based TABL model.

\begin{table}[h]
\centering
\begin{tabular}{lcccc}
\toprule
\textbf{Model} & \textbf{Accuracy \%} & \textbf{Precision \%} & \textbf{Recall \%} & \textbf{F1 Score \%} \\
\midrule
SVM & - & 39.62 & 44.92 & 35.88 \\
MLP & - & 47.81 & 60.78 & 48.27 \\
LSTM & - & 60.77 & 75.92 & 66.33 \\
CNN & - & 56.00 & 45.00 & 44.00 \\
\textbf{DeepLOB} & \textbf{84.47} & \textbf{84.00} & \textbf{84.47\%} & \textbf{83.40\%} \\
\bottomrule
\end{tabular}
\end{table}

In Setup 2 (fixed train/test split), DeepLOB shows a \textbf{5.7\% gain} over TABL. Furthermore, naive trading simulations yield statistically significant profits ($t > 10$), confirming the model's economic value.

\subsection{Transfer Learning \& Speed}
\begin{itemize}
    \item \textbf{Transfer}: When trained on LSE stocks and tested on \textit{unseen} stocks, performance drops only \textbf{1.55\%}, suggesting the model learns universal market physics.
    \item \textbf{Speed}: Inference takes 0.25 ms/step on GPU, suitable for high-frequency trading.
\end{itemize}

\subsection{Why DeepLOB Works ?}
The synergistic combination of components is critical for performance. Removing any single component significantly degrades results (Setup 2, k=10 F1):
\begin{enumerate}
    \item \textbf{Convolutional Layer Importance}: Removing the CNN layers and feeding raw 40-dim features to LSTM results in \textbf{66.33\%} F1 (-17\%). The CNN effectively extracts spatial structure and reduces noise before temporal processing.
    \item \textbf{Temporal Layer Importance}: Removing the LSTM/Inception layers and using only CNNs results in \textbf{55.21\%} F1 (-28\%). Spatial features alone cannot capture the evolution of order book states.
    \item \textbf{Multi-Scale Importance}: Replacing the Inception module with standard fixed-window convolutions (or naive attention as in TABL) results in lower performance (\textbf{77.63\%} for TABL vs 83.40\% for DeepLOB). The Inception module's ability to \emph{simultaneously} see short-term microstructure and long-term trends is a key differentiator.
\end{enumerate}

\section{Future Research Directions}

Several promising extensions could address current limitations:

\subsection{Bayesian DeepLOB (BDLOB)}
DeepLOB currently provides point estimates without uncertainty. 
BDLOB would replace standard weight layers with Bayesian layers to learn distributions over weights:
\[ y = \mathbb{E}_w[f(x; w)], \quad w \sim p(w|\mathcal{D}) \]
This allows the model to output confidence intervals, enabling \textbf{risk-aware position sizing} (scaling down bets when the model is uncertain).

\subsection{Attention Mechanisms}
While the Inception module uses fixed parallel kernels, an \textbf{attention mechanism} could dynamically weight different LOB levels or time steps based on the current market regime. 
For example, in high-volatility regimes, the deep levels (L5-L10) might be more informative than the best bid/ask.

\subsection{Hierarchical Temporal Modeling}
The current LSTM has a fixed 100-step window (~20s). A hierarchical architecture could model multiple timescales explicitly:
\begin{itemize}
    \item \textbf{Micro-LSTM}: Captures tick-by-tick microstructure.
    \item \textbf{Macro-LSTM}: Aggregates micro-states to capture minute-level trends or hourly regimes.
\end{itemize}

\subsection{Online Learning}
Financial markets are non-stationary. A static model trained on 2010 data will fail in 2024. 
Implementing \textbf{online learning} or periodic re-training pipelines is essential for production deployment to adapt to shifting market regimes.

\section{Strengths and Weaknesses}

\textbf{Strengths}:
\begin{itemize}
    \item \textbf{Accuracy}: State-of-the-art on FI-2010.
    \item \textbf{Efficiency}: Low parameter count (60k) reduces overfitting.
    \item \textbf{Transferability}: Works well on unseen assets.
    \item \textbf{Design}: Architecture reflects domain knowledge (stride, timescales).
\end{itemize}

\textbf{Weaknesses}:
\begin{itemize}
    \item \textbf{Black Box}: Lacks intrinsic interpretability.
    \item \textbf{Stationarity}: No explicit regime-switching mechanism.
    \item \textbf{Sensitivity}: Requires careful tuning of hyperparameters.
\end{itemize}



\section{Conclusion}
DeepLOB succeeds by replacing generic deep learning stacks with domain-aware components. 
The specific use of stride for feature pairing, Inception for timescale aggregation, and LSTM for efficiency creates a robust model for LOB prediction.

\section*{References}

\textbf{Core Paper}: Zhang, Z., Zohren, S., \& Roberts, S. (2020). DeepLOB: Deep Convolutional Neural Networks for Limit Order Books. \textit{arXiv:1808.03668v6}.

\textbf{Key Related Work}:
\begin{itemize}
    \item Ntakaris, A., et al. (2018). Benchmark Dataset for Mid-Price Forecasting of Limit Order Book Data.
    \item Szegedy, C., et al. (2014). Going Deeper with Convolutions (Inception Module).
    \item Hochreiter, S., \& Schmidhuber, J. (1997). Long Short-Term Memory.
    \item Ribeiro, M. T., et al. (2016). ``Why Should I Trust You?'': Explaining Predictions (LIME).
\end{itemize}

\end{document}
